\subsection{Convolution and linear representations}

PROPOSITION 11. --- Let G be a locally compact group, E a quasi-complete locally convex space, U a continuous representation of G in E.

\begin{enumerate}
\def\labelenumi{(\roman{enumi})}
\item
  If \(\lambda \in \mathcal{C}'(\mathbf{G})\) , \(\mu \in \mathcal{C}'(\mathbf{G})\) , then \(U(\lambda * \mu) = U(\lambda)U(\mu)\) .
\item
  Suppose that \(\mathbf{E}\) is a Banach space, and let \(\rho(s) = \|U(s)\|\) for \(s\in\mathbf{G}\) . If \(\lambda\in\mathcal{M}^{\rho}(\mathbf{G})\) , \(\mu\in\mathcal{M}^{\rho}(\mathbf{G})\) , then \(U(\lambda*\mu)=U(\lambda)U(\mu)\) .
\end{enumerate}

Let \(\lambda, \mu\) be in \(\mathcal{C}'(\mathrm{G})\) . For any \(x \in \mathrm{E}\) one has, by applying notably Props. 1 and 4 of Ch. VI, §1, No.~1,

\[
\begin{array}{r l} & U (\lambda * \mu) x = \int_ {\mathbf {G}} U (s) x d (\lambda * \mu) (s) \\ & \quad = \int_ {\mathbf {G} \times \mathbf {G}} U (s t) x d \lambda (s) d \mu (t) = \int_ {\mathbf {G} \times \mathbf {G}} U (s) U (t) x d \lambda (s) d \mu (t) \\ & \quad = \int_ {\mathbf {G}} U (\lambda) U (t) x d \mu (t) = U (\lambda) \int_ {\mathbf {G}} U (t) x d \mu (t) = U (\lambda) U (\mu) x, \end{array}
\]

whence (i). An analogous argument may be applied in case (ii).

With G still a locally compact group, let us assume that G operates continuously on the left in a locally compact space X. This defines (§2, No.~4) a continuous linear representation \(\pmb{\gamma}\) of G in \(\mathcal{M}(\mathrm{X})\) (equipped with the topology of compact convergence in \(\mathcal{K}(\mathrm{X})\) ).

PROPOSITION 12. --- If \(\lambda \in \mathcal{C}'(\mathbf{G})\) and \(\mu \in \mathcal{M}(\mathbf{X})\) , then

\[
\gamma (\lambda) \mu = \lambda * \mu .
\]

By Prop. 7 of §1, No.~5,

\[
\lambda * \mu = \int_ {G} (\varepsilon_ {s} * \mu) d \lambda (s).
\]

Now, \(\varepsilon_{s}*\mu = \pmb {\gamma}(s)\mu\) (No.~2, formula (5)) and

\[
\int_ {\mathsf {G}} \left(\pmb {\gamma} (s) \mu\right) d \lambda (s) = \pmb {\gamma} (\lambda) \mu
\]

by the definition of \(\pmb{\gamma}(\lambda)\) .

COROLLARY. --- The mapping \((\lambda,\mu)\mapsto\lambda*\mu\) of \(\mathcal{C}^{\prime}(\mathrm{G})\times\mathcal{M}(\mathrm{X})\) into \(\mathcal{M}(\mathrm{X})\) is hypocontinuous relative to the equicontinuous subsets of \(\mathcal{C}^{\prime}(\mathrm{G})\) and the compact subsets of \(\mathcal{M}(\mathrm{X})\) ( \(\mathcal{C}^{\prime}(\mathrm{G})\) and \(\mathcal{M}(\mathrm{X})\) being equipped with the topology of compact convergence in \(\mathcal{C}(\mathrm{G})\) and \(\mathcal{H}(\mathrm{X})\) , respectively).

For, \(\mathcal{M}(\mathrm{X})\) , equipped with the topology of compact convergence in \(\mathcal{K}(\mathrm{X})\) , is quasi-complete. Therefore the mapping \((\lambda, \mu) \mapsto \pmb{\gamma}(\lambda)\mu\) of \(\mathcal{C}'(\mathrm{G}) \times \mathcal{M}(\mathrm{X})\) into \(\mathcal{M}(\mathrm{X})\) is hypocontinuous relative to the equicontinuous subsets of \(\mathcal{C}'(\mathrm{G})\) and the compact subsets of \(\mathcal{M}(\mathrm{X})\) (§2, No.~6). It then suffices to apply Prop. 12.

Remarks. --- 1) Let \(\lambda_0 \in \mathcal{C}'(\mathbf{G})\) . The mapping \(\mu \mapsto \lambda_0 * \mu\) of \(\mathcal{M}(\mathbf{X})\) into \(\mathcal{M}(\mathbf{X})\) is vaguely continuous. For, let \(f \in \mathcal{K}(\mathbf{X})\) . One has

\[
\langle \lambda_ {0} * \mu , f \rangle = \int f (s x) d \lambda_ {0} (s) d \mu (x) = \langle \mu , g \rangle ,
\]

where \(g(x) = \int f(sx) d\lambda_0(s)\) . Now, \(g\) is continuous (Ch. VII, §1, No.~1, Lemma 1). On the other hand, let \(S\) be the support of \(\lambda_0\) and \(K\) that of \(f\) . The conditions \(sx \in K\) and \(s \in S\) imply \(x \in S^{-1}K\) ; therefore the support of \(g\) is contained in \(S^{-1}K\) , so that \(g \in \mathcal{K}(X)\) . Then \(\langle \lambda_0 * \mu, f \rangle = \langle \mu, g \rangle\) is a vaguely continuous function of \(\mu\) , which proves our assertion.

\begin{enumerate}
\def\labelenumi{\arabic{enumi})}
\setcounter{enumi}{1}
\tightlist
\item
  Let \(\mu_0 \in \mathcal{M}(\mathrm{X})\) . The mapping \(\lambda \mapsto \lambda * \mu_0\) of \(\mathcal{C}'(\mathrm{G})\) into \(\mathcal{M}(\mathrm{X})\) is continuous for the topologies \(\sigma(\mathcal{C}'(\mathrm{G}), \mathcal{C}(\mathrm{G}))\) and \(\sigma(\mathcal{M}(\mathrm{X}), \mathcal{K}(\mathrm{X}))\) . For, let \(f \in \mathcal{K}(\mathrm{X})\) . Setting \(h(s) = \int f(sx) d\mu_0(x)\) , we have \(\langle f, \lambda * \mu_0 \rangle = \langle h, \lambda \rangle\) , and \(h \in \mathcal{C}(\mathrm{G})\) (Ch. VII, §1, No.~1, Lemma 1).
\end{enumerate}

PROPOSITION 13. --- The mapping \((s,\mu)\mapsto\boldsymbol{\gamma}(s)\mu\) of \(\mathbf{G}\times\mathcal{M}_{+}(\mathbf{X})\) into \(\mathcal{M}_{+}(\mathbf{X})\) is continuous when the set \(\mathcal{M}_{+}(\mathbf{X})\) of positive measures on X is equipped with the vague topology.

Since \(\pmb{\gamma}(s)\mu = \pmb{\gamma}(ss_0^{-1})\pmb{\gamma}(s_0)\mu\) , it follows from Remark 1 that it suffices to prove the continuity of the mapping under consideration at a point of the form \((e,\mu_0)\) with \(\mu_0\in \mathcal{M}_+(\mathrm{X})\) . Given a function \(f\in \mathcal{K}(\mathrm{X})\) and a number \(\varepsilon >0\) , it is thus a matter of showing that there exist a neighborhood U of \(e\) in G and a neighborhood W of \(\mu_0\) in \(\mathcal{M}_+(\mathrm{X})\) such that the relations \(s\in \mathrm{U}\) , \(\mu \in \mathrm{W}\) imply

\[
\left| \int f (s x) d \mu (x) - \int f (x) d \mu_ {0} (x) \right| \leqslant \varepsilon .\tag{6}
\]

Let \(\mathbf{V}\) be a compact neighborhood of the support \(\mathbf{K}\) of \(f\) in \(\mathbf{X}\) , and let \(\varphi \in \mathcal{K}_{+}(\mathbf{X})\) be such that \(\varphi(x) = 1\) on \(\mathbf{V}\) ; there exists a neighborhood \(\mathrm{W}_0\) of \(\mu_0\) in \(\mathcal{M}_{+}(\mathbf{X})\) such that \(a = \sup_{\mu \in \mathrm{W}_0} \mu(\mathrm{V})\) is finite: it suffices to take for \(\mathrm{W}_0\) the set of \(\mu \in \mathcal{M}_{+}(\mathbf{X})\) such that \(|\langle \varphi, \mu - \mu_0 \rangle| \leqslant 1\) . Since the mapping \((s, x) \mapsto sx\) is continuous, there is, on the other hand, a compact neighborhood \(\mathrm{U}_0\) of \(e\) in \(\mathbf{G}\) such that \(s\mathrm{K} \subset \mathrm{V}\) for all \(s \in \mathrm{U}_0\) ; the function \((s, x) \mapsto f(sx)\) is then uniformly continuous in \(\mathrm{U}_0 \times \mathrm{V}\) and so there is a neighborhood \(\mathrm{U} \subset \mathrm{U}_0\) of \(e\) such that \(|f(sx) - f(x)| \leqslant \varepsilon / 2a\) for all \(s \in \mathrm{U}\) and \(x \in \mathrm{V}\) . For \(s \in \mathrm{U}\) and \(\mu \in \mathrm{W}_0\) , we therefore have

\[
\left| \int f (s x) d \mu (x) - \int f (x) d \mu (x) \right| \leqslant \varepsilon / 2;
\]

if \(\mathbf{W} \subset \mathbf{W}_0\) is the neighborhood of \(\mu_0\) in \(\mathcal{M}_{+}(\mathbf{X})\) formed by the measures \(\mu \in \mathbf{W}_0\) such that \(|\int f(x) d\mu(x) - \int f(x) d\mu_0(x)| \leqslant \varepsilon/2\) , \(\mathbf{U}\) and \(\mathbf{W}\) meet the requirements.

